Die Formelsammlung gibt die Dichte bei \TAmO, $\rho_0 = \rhMeO$ (FoTaBe S.~187). Bei dieser Temperatur ist das Volumen
\begin{align}
 V_0 &= \VnuF = \frac{\mMe}{\rhMe} = \Vnu
\end{align}
Beim Erwärmen um $\Delta\vartheta = \dTmF = \dTm$ dehnt es sich mit $\gamma = \gaMeO$ (FoTaBe S.~191) aus auf
\begin{align}
 V &= \VwaF \\
   &= \Vnu\cdot\left(1+\gaMe\cdot\dTm\right) \\
   &= \Vwa \approx \result{\VwacP{0}{2}}
\end{align}
